In French residential markets, listing prices anchor negotiation rather than serving as posted take-it-or-leave-it offers, and the transaction price conditional on a sale typically falls below the listing price by a margin that depends on local market conditions and property characteristics.\footnote{The negotiation margin in France is empirically small. The majority of properties transact at or very near list price, and the average discount rarely exceeds three to five percent. \citet{merlo_bargaining_2004} document similar patterns for England, showing that most residential transactions involve little or no negotiation once a listing price is set. \citet{genesove_loss_2001} provide evidence that list prices act as anchors that constrain final transaction prices. For France specifically, \citet{friggit_marche_2013} documents that the average negotiation margin has remained narrow and declining, with a large share of properties selling without any price reduction.} Buyer arrival and match quality are stochastic, so the seller faces uncertainty over both whether a sale occurs and, conditional on a sale, at what price. Her problem therefore separates into a transaction probability $\Gamma$ and an expected transaction price $\phi$, each depending on the listing price and disclosure decision:
